\begingroup
\setlength{\abovedisplayskip}{4pt plus 1pt minus 1pt}
\setlength{\belowdisplayskip}{4pt plus 1pt minus 1pt}
\setlength{\abovedisplayshortskip}{2pt plus 1pt}
\setlength{\belowdisplayshortskip}{2pt plus 1pt}
\setlength{\parskip}{0pt}
\section{\HMTLang{Memoria bilateral, holonomía relativa y restitución orientada}{Bilateral memory, relative holonomy, and oriented recovery}}
\label{delta22:xi:holonomia}
\HMTLang{La conservación bilateral adquiere contenido operacional cuando se conserva, además del estado visible, el canal que registra la composición. La selección excepcional y temporal de \(K\), desarrollada anteriormente, determina el registro; sus lectores aritméticos e incidenciales transmiten relaciones distintas del mismo estado enriquecido. Aquí se explicita cómo un orden de operaciones deja una respuesta recuperable incluso cuando coinciden sus multiplicidades y retorna la fase.}{Bilateral conservation acquires operational content when the channel recording composition is retained alongside the visible state. The exceptional and temporal selection of \(K\), developed above, determines the record; its arithmetic and incidence readers transmit distinct relations of the same enriched state. Here we make explicit how an ordering of operations leaves a recoverable response even when multiplicities agree and phase returns.}
\HMTLang{El antecedente es la generación APP--TRIT--TPK, con sus dos proyecciones, la prolongación completa y la orientación dodecafásica. La holonomía nonádica \(\Gamma_9\) pertenece a esa dinámica enriquecida. El conmutador relativo que se construye en este apartado corresponde a una realización operatoria especificada; su dominio, su forma conservada y su lector se mantienen diferenciados.}{The antecedent is APP--TRIT--TPK generation, with both projections, the complete prolongation, and dodecaphasic orientation. Nonadic holonomy \(\Gamma_9\) belongs to that enriched dynamics. The relative commutator constructed in this section belongs to a specified operator realization; its domain, conserved form, and reader remain distinguished.}

\subsection{\HMTLang{Composición unitaria con archivo bilateral}{Unitary composition with bilateral archive}}
\HMTLang{Para un transporte unitario \(C\) en un espacio de Hilbert \(\mathcal H\), sean}{For a unitary transport \(C\) on a Hilbert space \(\mathcal H\), let}
\[
 T_C=\frac{8I+C}{9},\qquad D_C=\frac{\sqrt8(C-I)}9,
 \qquad R_C=\frac{I+8C}{9},\qquad
 \mathbb U_C=\begin{pmatrix}T_C&D_C\\D_C&R_C\end{pmatrix}.
\]
\begin{theorem}[\HMTLang{Representación bilateral y composición completa}{Bilateral representation and complete composition}]
\HMTLang{El operador \(\mathbb U_C\) es unitario y satisface}{The operator \(\mathbb U_C\) is unitary and satisfies}
\[
 \mathbb U_{C_2}\mathbb U_{C_1}=\mathbb U_{C_2C_1},\qquad
 T_{C_2}T_{C_1}+D_{C_2}D_{C_1}=\frac{8I+C_2C_1}{9}.
\]
\HMTLang{Para un estado con archivo inicialmente nulo,}{For a state with initially zero archive,}
\[
 \|T_Cv\|^2+\|D_Cv\|^2=\|v\|^2,\qquad
 v=T_C^*T_Cv+D_C^*D_Cv.
\]
\end{theorem}
\begin{proof}
\HMTLang{Las matrices}{The matrices}
\[
 P_0=\frac19\begin{pmatrix}8&-\sqrt8\\-\sqrt8&1\end{pmatrix},
 \qquad P_1=\frac19\begin{pmatrix}1&\sqrt8\\\sqrt8&8\end{pmatrix}
\]
\HMTLang{son proyectores ortogonales complementarios: la multiplicación verifica \(P_i^2=P_i\), \(P_0P_1=0\), \(P_0+P_1=I_2\). Puesto que \(\mathbb U_C=P_0\otimes I+P_1\otimes C\), su producto adjunto es la identidad y la composición multiplica únicamente los factores \(C\). La primera componente del producto da la fórmula bilateral. Aplicar la unitariedad a \((v,0)\) prueba la igualdad de normas; su bloque superior izquierdo es \(T_C^*T_C+D_C^*D_C=I\), que da el recuperador.}{are complementary orthogonal projectors: multiplication verifies \(P_i^2=P_i\), \(P_0P_1=0\), and \(P_0+P_1=I_2\). Since \(\mathbb U_C=P_0\otimes I+P_1\otimes C\), its adjoint product is the identity, and composition multiplies only the factors \(C\). The first component of the product gives the bilateral formula. Applying unitarity to \((v,0)\) proves norm conservation; its upper-left block is \(T_C^*T_C+D_C^*D_C=I\), which gives the recovery map.}
\end{proof}
\HMTLang{El orden de composición queda distribuido de manera exacta:}{Composition order is distributed exactly:}
\[
 [T_{C_2},T_{C_1}]=\frac{[C_2,C_1]}{81},\qquad
 [D_{C_2},D_{C_1}]=\frac{8[C_2,C_1]}{81}.
\]
\HMTLang{Al sumar ambas contribuciones se obtiene \([C_2,C_1]/9\). Los factores \(1\) y \(8\) comparan estos términos operatorios; sus interpretaciones energéticas dependen de la forma y del lector. El archivo conserva la contribución que una composición de contracciones visibles por sí sola omitiría.}{Adding both contributions gives \([C_2,C_1]/9\). The factors \(1\) and \(8\) compare these operator terms; their energy interpretations depend on the form and reader. The archive retains the contribution that composition of visible contractions alone would omit.}

\subsection{\HMTLang{Plano de cuarto de giro en el registro dodecafásico}{Quarter-turn plane in the dodecaphasic record}}
\HMTLang{Sea \(C_{12}e_j=e_{j+1}\), con índices módulo \(12\), la permutación del calendario. El proyector sobre el período cuatro y su componente antiperiódica son}{Let \(C_{12}e_j=e_{j+1}\), with indices modulo \(12\), be the calendar permutation. The projector onto period four and its antiperiodic component are}
\[
 P_4=\frac{I+C_{12}^4+C_{12}^8}{3},\qquad
 Q=P_4\frac{I-C_{12}^6}{2}.
\]
\HMTLang{Definamos}{Define}
\[
 u_0=(1,0,-1,0,1,0,-1,0,1,0,-1,0)^{\mathsf T},
 \quad v_0=C_{12}u_0,\quad W=\frac1{\sqrt6}(u_0,v_0).
\]
\begin{proposition}[\HMTLang{Realización exacta del cuarto de giro}{Exact quarter-turn realization}]
\[
 W^{\mathsf T}W=I_2,\quad WW^{\mathsf T}=Q,\quad
 C_{12}W=WJ_0,\quad C_{12}^3W=-WJ_0,\qquad
 J_0=\begin{pmatrix}0&-1\\1&0\end{pmatrix}.
\]
\end{proposition}
\begin{proof}
\HMTLang{Los vectores \(u_0,v_0\) tienen seis entradas de módulo uno y soportes disjuntos, por lo que son ortogonales y tienen norma \(\sqrt6\). Además \(C_{12}v_0=-u_0\). El rango de \(Q\) contiene precisamente los vectores de período cuatro \((a,b,-a,-b)\) repetidos tres veces; por tanto tiene dimensión dos y coincide con la imagen de \(W\). Las identidades restantes se deducen por aplicación a sus dos columnas.}{The vectors \(u_0,v_0\) each have six entries of modulus one and disjoint supports, so they are orthogonal with norm \(\sqrt6\). Moreover, \(C_{12}v_0=-u_0\). The range of \(Q\) consists precisely of period-four vectors \((a,b,-a,-b)\) repeated three times; it is therefore two-dimensional and equals the image of \(W\). The remaining identities follow by applying the operators to both columns.}
\end{proof}
% Fragmento ES para inserción, no documento autónomo.
% Destino: XII, ampliacion_20260922.tex, después de la prueba de
% «Realización exacta del cuarto de giro» y antes de introducir S_zeta.
% Propietarios y comprobaciones: LEER_PRIMERO.md.
\subsubsection{Restitución del calendario y acarreo ternario}
\label{ins29:xii:acarreo}

La extracción del plano de cuarto de giro admite una descripción exacta
de la información temporal conservada. Recibimos el calendario de
\(108\) posiciones y la extensión dodecafásica sobre la fase nonádica
construidos anteriormente. Para esta descripción utilizamos representantes
de base cero:
\[
 \theta=r+9\beta,\qquad 0\leq r<9,\quad 0\leq\beta<12.
\]
La posición que el registro enumera de \(1\) a \(9\) es aquí \(r+1\).
La actualización enriquecida precede a este calendario finito y conserva,
además, la ruta, las incidencias y la memoria acumulada. En el plano ya
construido, \(C_{12}^{\beta}W=WJ_0^{\beta}\); por tanto, su lectura depende
de \(b=\beta\bmod4\).

\begin{proposition}[Fibra ternaria de la lectura de cuarto de giro]
\label{ins29:xii:fibra}
El par formado por la fase nonádica \(r\) y la fase planar \(b\) determina
un cociente de \(36\) posiciones del calendario. Su ley de composición es
\[
 (r,b)(r',b')=
 \left(r+r'\bmod9,\;
 b+b'+\left\lfloor\frac{r+r'}9\right\rfloor\bmod4\right).
\]
Cada lectura posee tres preimágenes en \(\mathbb Z/108\mathbb Z\).
La componente adicional \(k=\beta\bmod3\) restituye la coordenada
dodecafásica mediante
\[
 \beta=9b+4k\pmod{12}.
\]
\end{proposition}
\begin{proof}
La suma de dos representantes da
\(\theta+\theta'=r+r'+9(\beta+\beta')\). La división euclidiana de
\(r+r'\) por \(9\) produce el acarreo de la fórmula. La lectura
\((r,b)\) equivale al residuo \(\theta\bmod36\), pues
\(\theta\equiv r+9b\pmod{36}\). Su núcleo es
\(\{0,36,72\}\), de modo que cada fibra tiene tres elementos.
Finalmente, \(9b+4k\) es congruente con \(b\) módulo \(4\) y con
\(k\) módulo \(3\); estos dos residuos determinan un único elemento
módulo \(12\).
\end{proof}

\begin{proposition}[Localización de la extensión no escindida]
\label{ins29:xii:extension}
La descomposición
\[
 \theta\longmapsto(j,z)=(\theta\bmod4,\theta\bmod27),
 \qquad \theta=81j+28z\pmod{108},
\]
identifica el calendario con
\(\mathbb Z/4\mathbb Z\times\mathbb Z/27\mathbb Z\).
La proyección nonádica es \((j,z)\mapsto z\bmod9\), y la inclusión
dodecafásica \(\beta\mapsto9\beta\) toma la forma
\[
 \beta\longmapsto(\beta\bmod4,9\beta\bmod27).
\]
La obstrucción a escindir la extensión completa queda concentrada en
\[
 0\longrightarrow\mathbb Z/3\mathbb Z
 \longrightarrow\mathbb Z/27\mathbb Z
 \longrightarrow\mathbb Z/9\mathbb Z\longrightarrow0.
\]
\end{proposition}
\begin{proof}
Los coeficientes \(81\) y \(28\) tienen residuos \((1,0)\) y \((0,1)\)
módulo \((4,27)\), respectivamente. Esto verifica la inversa y las dos
aplicaciones indicadas. La sucesión ternaria no se escinde: un producto
\(\mathbb Z/3\mathbb Z\times\mathbb Z/9\mathbb Z\) tiene exponente
\(9\), mientras \(\mathbb Z/27\mathbb Z\) tiene exponente \(27\).
En cambio, la proyección del cociente de \(36\) posiciones sobre
\(\mathbb Z/9\mathbb Z\) admite la sección homomorfa
\(n\mapsto28n\pmod{36}\). Por tanto, la reducción elimina precisamente
esta obstrucción ternaria.
\end{proof}

La distinción se observa ya en \(\theta=0,36,72\): las tres posiciones
dan \(r=b=0\), mientras \(z\) vale \(0,9,18\), respectivamente.
La restitución del calendario requiere conservar la fibra que el plano
identifica. Además, el residuo cuaternario global satisface
\(j=r+b\pmod4\); la fase planar \(b\) se transporta junto con la fase
nonádica y su acarreo.

Así queda localizado el alcance del cuarto de giro: organiza una
representación orientada del registro y, acompañado del residuo nonádico,
recupera exactamente su cociente de \(36\) posiciones. La coordenada
ternaria adicional restituye las \(108\) posiciones. La memoria de la
dinámica enriquecida conserva a su vez la prolongación entre retornos
completos del calendario. La periodicidad de una representación y la
persistencia de la trayectoria registrada pertenecen, por ello, a lecturas distintas
del mismo transporte.

\HMTLang{El plano se extrae de la permutación ya definida. Su forma de acción se recibe del par angular generado, con \(A\) y \(C^*\) expresados en la misma unidad, \(|C^*/A|<1\) y \(\theta\in\mathbb R\):}{The plane is extracted from the already defined permutation. Its action form is received from the generated angular pair, with \(A\) and \(C^*\) expressed in the same unit, \(|C^*/A|<1\), and \(\theta\in\mathbb R\):}
\[
 \xi=C^*/A,\quad \zeta=\operatorname{artanh}\xi,\quad
 S_\zeta=\operatorname{diag}(e^{\zeta/2},e^{-\zeta/2}),\quad
 U_\zeta(\theta)=S_\zeta e^{-\theta J_0}S_\zeta^{-1}
 =\begin{pmatrix}\cos\theta&e^\zeta\sin\theta\\
 -e^{-\zeta}\sin\theta&\cos\theta\end{pmatrix}.
\]
\HMTLang{Cada \(U_\zeta\) conserva \(g_\zeta=S_\zeta^{-\mathsf T}S_\zeta^{-1}\) y tiene determinante uno. Su forma conservada queda especificada antes de comparar transportes con distinta anisotropía.}{Each \(U_\zeta\) preserves \(g_\zeta=S_\zeta^{-\mathsf T}S_\zeta^{-1}\) and has determinant one. Its conserved form is specified before comparing transports with different anisotropy.}

\HMTLang{Los productos operatorios actúan de derecha a izquierda.}{Operator products act from right to left.}

\subsection{\HMTLang{Conmutador relativo y restitución de la acción}{Relative commutator and action recovery}}
\begin{theorem}[\HMTLang{Defecto de conmutación de las formas conjugadas}{Commutation defect of conjugate forms}]
\HMTLang{Sean \(U=U_{\zeta_s}(\theta)\), \(V=U_{\zeta_r}(\phi)\), \(H=UVU^{-1}V^{-1}\). Con \(\delta=\zeta_s-\zeta_r\) y \(\chi=\sinh\delta\sin\theta\sin\phi\),}{Let \(U=U_{\zeta_s}(\theta)\), \(V=U_{\zeta_r}(\phi)\), and \(H=UVU^{-1}V^{-1}\). With \(\delta=\zeta_s-\zeta_r\) and \(\chi=\sinh\delta\sin\theta\sin\phi\),}
\[
 \det H=1,\qquad \operatorname{tr}H=2+4\chi^2,\qquad
 \lambda_\pm(H)=(\sqrt{1+\chi^2}\pm|\chi|)^2.
\]
\HMTLang{Además, \(H=I\) exactamente cuando \(\chi=0\).}{Moreover, \(H=I\) precisely when \(\chi=0\).}
\end{theorem}
\begin{proof}
\HMTLang{La multiplicación de las matrices da \(\operatorname{tr}(UV)=2(\cos\theta\cos\phi-\cosh\delta\sin\theta\sin\phi)\). Para matrices de determinante uno, Cayley--Hamilton da \(A^{-1}=(\operatorname{tr}A)I-A\). Sustituir esta expresión para las dos inversas produce}{Multiplication gives \(\operatorname{tr}(UV)=2(\cos\theta\cos\phi-\cosh\delta\sin\theta\sin\phi)\). For determinant-one matrices, Cayley--Hamilton gives \(A^{-1}=(\operatorname{tr}A)I-A\). Substituting this expression for both inverses yields}
\[
 \operatorname{tr}(UVU^{-1}V^{-1})
 =a^2+b^2+c^2-abc-2,
 \quad a=\operatorname{tr}U,\ b=\operatorname{tr}V,\ c=\operatorname{tr}(UV).
\]
\HMTLang{La sustitución trigonométrica reduce la expresión a \(2+4\chi^2\). Resolver \(\lambda^2-(2+4\chi^2)\lambda+1=0\) da los autovalores. Si \(\chi=0\), o coinciden las formas, o uno de los transportes es \(\pm I\); ambos casos conmutan. La implicación inversa se deduce de la traza.}{Trigonometric substitution reduces this expression to \(2+4\chi^2\). Solving \(\lambda^2-(2+4\chi^2)\lambda+1=0\) gives the eigenvalues. If \(\chi=0\), either the forms coincide or one transport is \(\pm I\); both cases commute. The converse follows from the trace.}
\end{proof}
\HMTLang{Para los cuartos de giro conjugados \(U_+=U_\zeta(\pi/2)\), \(U_-=U_{-\zeta}(\pi/2)\), las inversas son \(-U_\pm\). Por multiplicación,}{For conjugate quarter turns \(U_+=U_\zeta(\pi/2)\), \(U_-=U_{-\zeta}(\pi/2)\), the inverses are \(-U_\pm\). Multiplication gives}
\begin{equation}
 H=(U_+U_-)^2=\operatorname{diag}(e^{4\zeta},e^{-4\zeta}),\qquad
 \operatorname{tr}H-2=\frac{16\xi^2}{(1-\xi^2)^2}.
 \label{delta22:xi:lazo}
\end{equation}
\HMTLang{Con las representaciones \(A=1000\alpha\), \(C^*=2(\eta_{\rm ret}+\alpha)\), ambas en grados, resulta}{With the outputs \(A=1000\alpha\), \(C^*=2(\eta_{\rm ret}+\alpha)\), both in degrees, this gives}
\begin{equation}
 \rho_{\rm lazo}=e^{4\zeta}
 =\left(\frac{501\alpha+\eta_{\rm ret}}{499\alpha-\eta_{\rm ret}}\right)^2,
 \qquad 0<\eta_{\rm ret}<499\alpha.
 \label{delta22:xi:rho}
\end{equation}
\HMTLang{Su inversa, con el eje orientado conservado, es}{Its inverse, retaining the oriented axis, is}
\[
 \eta_{\rm ret}=\alpha\left[
 500\frac{\sqrt{\rho_{\rm lazo}}-1}{\sqrt{\rho_{\rm lazo}}+1}-1\right],
 \qquad \hbar_{\rm ret}=s_0\eta_{\rm ret}.
\]
\HMTLang{La prueba usa \(e^{2\zeta}=(1+\xi)/(1-\xi)\) y despeja \(\xi\). Se conserva \(s_0=10^{-34}\mathcal U_S\), previamente construido. Bajo conjugación se transporta la etiqueta del eje; la traza sola conserva \(|\zeta|\). La razón \(\rho_{\rm lazo}\) designa esta respuesta relativa, mientras \(\pi/\varphi^2\) conserva la definición distinta de razón regional marcada.}{The proof uses \(e^{2\zeta}=(1+\xi)/(1-\xi)\) and solves for \(\xi\). The previously constructed \(s_0=10^{-34}\mathcal U_S\) is retained. Under conjugation, the axis label is transported; the trace alone retains \(|\zeta|\). The ratio \(\rho_{\rm lazo}\) denotes this relative response, while \(\pi/\varphi^2\) retains its distinct definition as the marked regional ratio.}

\HMTLang{La rama positiva corresponde exactamente a \(\rho_{\rm lazo}>(501/499)^2\). Al invertir la orientación, \(\rho_{\rm lazo}\mapsto\rho_{\rm lazo}^{-1}\), se transporta la etiqueta del eje antes de aplicar el recuperador de esa rama.}{The positive branch corresponds exactly to \(\rho_{\rm lazo}>(501/499)^2\). Under orientation reversal, \(\rho_{\rm lazo}\mapsto\rho_{\rm lazo}^{-1}\), the axis label is transported before applying the recovery map of that branch.}

\subsection{\HMTLang{Levantamiento dodecafásico y amplificación compatible}{Dodecaphasic lift and compatible amplification}}
\HMTLang{Definamos sobre las doce coordenadas}{Define on the twelve coordinates}
\[
 \widetilde S_\zeta=I-Q+WS_\zeta W^{\mathsf T},\qquad
 L_\pm=\widetilde S_{\pm\zeta}C_{12}^3\widetilde S_{\pm\zeta}^{-1}.
\]
\begin{proposition}[\HMTLang{Restitución del conmutador en el registro completo}{Restoration of the commutator in the full record}]
\[
 L_+L_-L_+^{-1}L_-^{-1}=I-Q+WHW^{\mathsf T}.
\]
\end{proposition}
\begin{proof}
\HMTLang{Los operadores \(Q\) y \(C_{12}\) conmutan. Sobre \(\operatorname{im}Q\), la isometría \(W\) identifica \(C_{12}^3\) con \(-J_0\), y \(L_\pm\) con \(U_\pm\). Sobre el complemento, ambos \(L_\pm\) coinciden con \(C_{12}^3\), de modo que su conmutador es la identidad. La suma ortogonal de ambas restricciones da la fórmula.}{The operators \(Q\) and \(C_{12}\) commute. On \(\operatorname{im}Q\), the isometry \(W\) identifies \(C_{12}^3\) with \(-J_0\), and \(L_\pm\) with \(U_\pm\). On the complement, both \(L_\pm\) equal \(C_{12}^3\), so their commutator is the identity. The orthogonal sum of both restrictions gives the formula.}
\end{proof}
\HMTLang{Para la amplificación declarada \(\mathcal H_m=\mathbb R^{12}\otimes(\mathbb R^9)^{\otimes m}\), sean}{For the specified amplification \(\mathcal H_m=\mathbb R^{12}\otimes(\mathbb R^9)^{\otimes m}\), let}
\[
 \widetilde H=I-Q+WHW^{\mathsf T},\qquad
 H_m=\widetilde H\otimes I_{9^m},\qquad
 J_m v=v\otimes u_9,\quad u_9=(1,\ldots,1)/3.
\]
\HMTLang{Entonces \(J_m\) es isométrico y \(H_{m+1}J_m=J_mH_m\). La norma del operador permanece fija al amplificar por una identidad; por tanto la acción compatible se extiende al límite de Hilbert. Las pruebas anteriores se conservan en este refinamiento tensorial, con los proyectores y el lector transportados.}{Then \(J_m\) is isometric and \(H_{m+1}J_m=J_mH_m\). Operator norm remains fixed under amplification by an identity, so the compatible action extends to the Hilbert limit. The preceding proofs are retained in this tensor refinement, with projectors and reader transported.}

\subsection{\HMTLang{Información cronológica y lectura energética}{Chronological information and energy reading}}
\HMTLang{La potencia del lazo es \(H^n=\operatorname{diag}(\rho_{\rm lazo}^n,\rho_{\rm lazo}^{-n})\). Conserva \(QP\) y \(dQ\wedge dP\). Para \(\rho_{\rm lazo}\ne1\), una preparación con componente \(Q_0\ne0\) y eje conocido permite recuperar \(n=\log(Q_n/Q_0)/\log\rho_{\rm lazo}\). La preparación y la escala del detector pertenecen al dominio de ese recuperador.}{The loop power is \(H^n=\operatorname{diag}(\rho_{\rm lazo}^n,\rho_{\rm lazo}^{-n})\). It preserves \(QP\) and \(dQ\wedge dP\). For \(\rho_{\rm lazo}\ne1\), a preparation with \(Q_0\ne0\) and a known axis recovers \(n=\log(Q_n/Q_0)/\log\rho_{\rm lazo}\). Preparation and detector scale belong to the domain of this recovery map.}

\HMTLang{Se admite \(n\in\mathbb Z\), incluyendo la iteración del lazo inverso. Cuando \(Q_0=0\) y \(P_0\ne0\), la otra componente proporciona}{Allow \(n\in\mathbb Z\), including iteration of the inverse loop. When \(Q_0=0\) and \(P_0\ne0\), the other component gives}
\[
 n=-\frac{\log|P_n/P_0|}{\log\rho_{\rm lazo}}.
\]

\HMTLang{Las dos ordenaciones \(UU^{-1}VV^{-1}\) y \(UVU^{-1}V^{-1}\) contienen las mismas cuatro multiplicidades, pero sus respuestas son \(I\) y \(H\). El archivo ordenado conserva una distinción que el recuento conmutativo elimina. Para un conjunto de detectores \(\mathcal D\), dos historias son operacionalmente equivalentes cuando todas sus lecturas en \(\mathcal D\) coinciden. Las respuestas sobre una base completa determinan el operador lineal; la recuperación de una historia enriquecida requiere además sus restantes registros.}{The two orderings \(UU^{-1}VV^{-1}\) and \(UVU^{-1}V^{-1}\) contain the same four multiplicities, but their responses are \(I\) and \(H\). The ordered archive retains a distinction removed by commutative counting. For a detector set \(\mathcal D\), two histories are operationally equivalent when all their readings in \(\mathcal D\) agree. Responses on a complete basis determine the linear operator; recovery of an enriched history additionally requires its remaining records.}

\HMTLang{La forma de referencia fija el significado energético. Para \(g_{\rm ref}=\operatorname{diag}(e^{-\zeta_r},e^{\zeta_r})\), las ganancias de energía cuadrática del lazo diagonal, sobre los ejes marcados, son}{The reference form fixes the energy meaning. For \(g_{\rm ref}=\operatorname{diag}(e^{-\zeta_r},e^{\zeta_r})\), the diagonal loop's quadratic-energy gains along the marked axes are}
\[
 G_Q=e^{8\zeta},\qquad G_P=e^{-8\zeta},\qquad
 \mathcal O=\tfrac14\log(G_Q/G_P)=4\zeta.
\]
\HMTLang{Estas ganancias son los cuadrados de las dilataciones de estado. \(H\) y \(H^{-1}\) tienen la misma traza, mientras \(\mathcal O\) conserva el signo orientado. Para \(\zeta\ne0\), \(U_+\) y \(U_-\) carecen de una métrica positiva común: si ambos la conservaran, su conmutador también la conservaría y todos sus autovalores tendrían módulo uno, en contradicción con \eqref{delta22:xi:lazo}. Cada transporte individual conserva su propia \(g_{\pm\zeta}\).}{These gains are the squares of the state dilations. \(H\) and \(H^{-1}\) have the same trace, whereas \(\mathcal O\) retains the oriented sign. For \(\zeta\ne0\), \(U_+\) and \(U_-\) have no common positive metric: if both preserved one, their commutator would also preserve it and all its eigenvalues would have modulus one, contradicting \eqref{delta22:xi:lazo}. Each individual transport preserves its own \(g_{\pm\zeta}\).}

\HMTLang{La identidad algebraica de composición bilateral sigue siendo válida para estos operadores. La conservación unitaria conjunta requiere, adicionalmente, una misma forma positiva conservada. La holonomía relativa conserva el área y sus propios datos de restitución.}{The algebraic identity of bilateral composition remains valid for these operators. Joint unitary conservation additionally requires a common preserved positive form. Relative holonomy preserves area and its own recovery data.}

\HMTLang{Un transporte pasivo \(T_n=S_{\zeta_{n+1}}R(-\theta_n)S_{\zeta_n}^{-1}\) telescopa: su producto es \(S_{\zeta_N}R(-\sum\theta_n)S_{\zeta_0}^{-1}\), igual a la identidad cuando retornan la forma y la fase total. El conmutador relativo compara, en cambio, operaciones efectivas sobre una referencia común. Atribuir a esa respuesta un intercambio energético exige incluir el sistema, la memoria y el dispositivo. La conservación del área y la reversibilidad operatoria quedan demostradas con independencia de esa atribución.}{A passive transport \(T_n=S_{\zeta_{n+1}}R(-\theta_n)S_{\zeta_n}^{-1}\) telescopes: its product is \(S_{\zeta_N}R(-\sum\theta_n)S_{\zeta_0}^{-1}\), equal to the identity when shape and total phase return. The relative commutator instead compares effective operations against a common reference. Attributing an energy exchange to that response requires including the system, memory, and device. Area conservation and operator reversibility are proved independently of that attribution.}

\HMTLang{El resultado reúne tres niveles de recuperación: el archivo bilateral completa la composición visible; el calendario dodecafásico determina un plano operativo exacto; el lector orientado del lazo recupera el contraste angular y la acción con sus antecedentes. Esta articulación prolonga la función de acoplamiento de \(K\) sin sustituir su selección por una inversión posterior. La distinción entre fase, orden de composición, forma conservada y memoria recuperable queda expresada mediante operadores y recuperadores explícitos.}{The result assembles three levels of recovery: the bilateral archive completes visible composition; the dodecaphasic calendar determines an exact operational plane; and the oriented loop reader recovers angular contrast and action with their antecedents. This articulation extends the coupling role of \(K\) without replacing its selection by a later inversion. The distinction between phase, composition order, conserved form, and recoverable memory is expressed by explicit operators and recovery maps.}
\par\endgroup
